\documentclass[12pt,letterpaper]{article}

% --- Paquetes ---
\usepackage[utf8]{inputenc}
\usepackage[T1]{fontenc}
\usepackage[spanish]{babel}
\usepackage[margin=2.5cm]{geometry}
\usepackage{float}
\usepackage{tabularx}
\usepackage{hyperref}
\usepackage{graphicx}
\usepackage{setspace}
\usepackage{parskip}
\usepackage{helvet}      % Fuente sans-serif para un look más limpio
\renewcommand{\familydefault}{\sfdefault}

% --- Datos de la portada ---
\newcommand{\institucion}{Universidad Nacional Autónoma de México}
\newcommand{\facultad}{Facultad de Ingeniería}
\newcommand{\numpractica}{Práctica 2}
\newcommand{\semestre}{Semestre 2027-1}
\newcommand{\entrega}{17 de septiembre de 2026}
\newcommand{\grupo}{Grupo: 16}
\newcommand{\alumno}{López Guevara Sebastián}
\newcommand{\profesora}{Sonia López Maldonado}
\newcommand{\practica}{Calderas y calorimetría}
\newcommand{\materia}{SISTEMAS DE CONVERSIÓN DE ENERGÍA TÉRMICA}

\begin{document}
\thispagestyle{empty}

\begin{center}

    % --- INSTITUCIÓN ---
    {\Large\bfseries \institucion}\\[0.3cm]
    {\Large\bfseries \facultad}\\[0.3cm]
    {\large \semestre}\\[0.2cm]
    {\large \grupo}\\[1.5cm]
    
    % --- LOGO (opcional) ---
    % \includegraphics[width=0.25\textwidth]{logo.png}\\[1.5cm]
    
    % --- TÍTULO DE LA PRÁCTICA ---
    {\large \numpractica}\\[0.2cm]
    {\Huge\bfseries \practica}\\[0.5cm]
    \rule{0.8\textwidth}{0.4pt}\\[1.5cm]

    % --- DATOS DEL ALUMNO Y PROFESORA ---
    \begin{tabular}{rl}
        \textbf{Alumno:}     & \alumno     \\[0.3cm]
        \textbf{Profesora:}  & \profesora  \\[0.3cm]
        \textbf{Materia:}    & \materia    \\
    \end{tabular}

    \vfill
    
    % --- FECHA ---
    {\entrega}
    
\end{center}

\break

% ============================
% OBJETIVO
% ============================
\section{Objetivo}
Conocer la capacidad nominal, capacidad real, factor de sobrecarga, factor de vaporización,
eficiencia y calidad del vapor
por medio de mediciones realizadas a un generador en operación y analizarlas para
poder afirmar la exactitud de tales mediciones con datos numéricos.

% MATERIAL
% ============================
\section{Material}
\begin{itemize}
    \item Termómetro de 150 °C
    \item Termómetro de 250 °C
    \item Flexómetro
    \item Cronómetro
    \item 3 Cubetas
\end{itemize}

% ============================
% DESARROLLO
% ============================
\section{Desarrollo}
\begin{enumerate}
    \item Se midió la temperatura de entrada del agua de alimentación.
    \item Se tomó como dato la superficie de calefacción de la caldera.
    \item Se colocó el termómetro de 150 grados en el calorímetro de estrangulamiento al costado de 
    la caldera. Se dejó estabilizarse y mientras se continuó al calorímetro eléctrico.
    \item Se encendió la resistencia del calorímetro eléctrico y se colocaron cubetas bajo la llave.
    \item Se colocó el termómetro de 250 grados en el calorímetro eléctrico.
    \item Se realizaron las lecturas del voltaje e intensidad 
    de corriente del calorímetro eléctrico.
    \item Se esperó a que se calentara el agua.
    \item Se realizó la medición de la temperatura del calorímetro de estrangulamiento.
    \item Una vez el agua ascendió hasta la temperatura $T_3$, se colocó una cubeta vacía y
    se tomó el tiempo con un cronómetro.
    \item Al pasar 30 segundos aproximadamente se retiró la cubeta y se pesó.
    \item Se vació la cubeta y se pesó nuevamente.
    \item Se tomaron las medidas del volumen para el gasto de combustible.
\end{enumerate}

\vfill

% ============================
% TABLA DE DATOS
% ============================
\section{Tabla de datos}

\begin{table}[H] \label{tab:data}
\centering
\renewcommand{\arraystretch}{1.4}
    \begin{tabular}{|l|l|}
    \hline
    \textbf{Medición} & \textbf{Variable} \\
    \hline
    $6.4\ [{kg/cm}^2]$   & $P_2$ \\
    \hline
    $46.61\ [{m}^2]$     & Sup. Calef \\
    \hline
    $32\ [^\circ{C}]$    & $T_1$ \\
    \hline
    $30\ [{cm}]$ & $a$ \\
    \hline
    $30\ [{cm}]$ & $a$ \\
    \hline
    $5\ [{cm}]$  & $A$ \\
    \hline
    $2\ [min] + 57\ [s]$  & $t$ \\
    \hline
    $120\ [^\circ{C}]$   & $T_3$ \\
    \hline
    $106\ [{V}]$         & $V$ \\
    \hline
    $15.5\ [{A}]$        & $I$ \\
    \hline
    $2.5\ [{kg/cm}^2]$   & $P_4$ \\
    \hline
    $159\ [^\circ{C}]$   & $T_4$ \\
    \hline
    $850\ [{g}]$         & $m_{vc}$ \\
    \hline
    $300\ [{g}]$         & $m_{cu}$ \\
    \hline
    $30\ [{s}]$          & $t$ \\
    \hline
    \end{tabular}
\caption{Datos experimentales obtenidos durante la práctica.}
\end{table}

% ============================
% CÁLCULOS
% ============================
\section{Cálculos}

\subsection{Obtención de las propiedades termodinámicas del estado I} \label{entalpia_I}
Conocemos la temperatura de entrada y por lo tanto realizamos:

\begin{equation} \label{eq:cp_equation}
    h_{1}=C_p (T_1 - 0)
\end{equation}

Sustituyendo la temperatura uno registrada

$$
    h_{1}=C_p (32[ ^\circ C] - 0)  
$$

Sabemos que el $C_p$ del agua tiene un valor de $4.186[\frac{kJ}{kg ^\circ C}]$. 
Resolviendo la ecuación \ref{eq:cp_equation}

$$
    h_1 = 4.186[\frac{kJ}{kg\cdot K}] \cdot (32+273.15-273.15)[K]
$$

$$
    h_1 = 133.952 [\frac{kJ}{kg}]
$$

\subsection{Obtención de las propiedades termodinámicas del estado II} \label{stateII}

En esta sección se encuentran los datos correspondientes al estado 2.
Para entrar a tablas se utiliza $P_2 = 6.4[\frac{kg}{cm^2}]$. 
Convertimos este a $P_2$ a $627.6256[kPa]$. 
Luego obtenemos la presión de la caldera absoluta 
$$
    P_{abs} = 78.51[kPa] + 627.6256[kPa] = 706.1356[kPa]
$$

Buscamos $h_f$ y $h_{fg}$ en tablas:
$$
    h_{g} = 2763.12 \left[\frac{kJ}{kg}\right]
$$
$$
    h_{f} = 698.53 \left[\frac{kJ}{kg}\right]
$$

\subsection{Obtención de las propiedades termodinámicas del estado III 
                y la calidad del estado II} \label{x_calorimetro_estrangulamiento}

Esta medición corresponde al calorímetro de estrangulamiento. Los datos que se usarán son

\begin{tabularx}{0.8\textwidth} { 
  | >{\raggedright\arraybackslash}X 
  | >{\centering\arraybackslash}X 
  | >{\raggedleft\arraybackslash}X | }
 \hline
 $P_3$ & $785.1[hPa]$  \\
 \hline
 $T_3$ & $120[^\circ C]$  \\
\hline
\end{tabularx}

La presión fue obtenida a partir del \href{http://189.254.33.151/stn/coyoacan/Current_Monitor.htm}{monitor meteorológico de Coyoacán}.

Convertimos esta a $[kPa]$ y obtenemos $78.51[kPa]$.

Entramos a tablas con estos valores para el estado 3 para obtener $h_3$:

$$
    h_3 = h_2 = 2718838.42 \left[\frac{J}{kg}\right] = 2718.83 \left[\frac{kJ}{kg}\right]
$$

Para encontrar la calidad del agua que se encuentra como mezcla dentro de la caldera en el estado 2, 
buscaremos a la calidad con ayuda de las propiedades del estado III.

Utilizaremos la ecuación 
\begin{equation} \label{eq:calidad}
    x = \frac{h_2-h_f}{h_{fg}}
\end{equation}

Nos apoyamos de la entalpía del estado 3, ya que el proceso 2-3 es un proceso
isentálpico, para sustituir en la ecuación \ref{eq:calidad}.

$$
    x = \frac{
        2718.83842 \left[\frac{kJ}{kg}\right] - 698.53 \left[\frac{kJ}{kg}\right] }{
            2763.12 \left[\frac{kJ}{kg}\right] - 698.53 \left[\frac{kJ}{kg}\right]
        } = 0.97855
$$

\subsection{Obtención de las propiedades termodinámicas del estado IV y 
            la calidad del estado II} \label{x_calorimetro_electrico}

Estas propiedades corresponden al estado 4. Las mediciones fueron tomadas del calorímetro
eléctrico.

Para poder sustituir dentro de la ecuación

\begin{equation} \label{eq:calorimetro_electrico}
    h_2 = h_4 - \frac{VI}{ \dot{M_{vc}}}    
\end{equation}

se usarán los datos

\begin{tabularx}{0.8\textwidth} { 
  | >{\raggedright\arraybackslash}X 
  | >{\centering\arraybackslash}X 
  | >{\raggedleft\arraybackslash}X | }
 \hline
 $V$ & $106[V]$  \\
 \hline
 $I$ & $15.5[A]$  \\
 \hline
 $M_{vc}$ & $850[g]$  \\
 \hline
 $M_{cu}$ & $300[g]$  \\
 \hline
 $t_{vc}$ & $30[s]$  \\
 \hline
 $P_4$ & $2.5\left[\frac{kg}{cm^2} \right]$  \\
 \hline
 $T_4$ & $159[^\circ C]$  \\
\hline
\end{tabularx}

Convertimos $[g]\to[kg]$ sustituimos las masas medidas y el tiempo para obtener flujo másico

$$
    \dot{M_{vc}}=\frac{M_{vc}-M_{cu}}{t_{vc}}=\frac{0.850[kg]-0.300[kg]}{30[s]}=0.01833 \left[\frac{kg}{s}\right]
$$

Buscamos la entalpía del estado 4. Para esto convertimos $P_4$ a $[kPa]$

$$
    P_4 = 2.5\left[\frac{kg}{cm^2} \right] = 245166.25[Pa] = 245.16625[kPa]
$$

y encontramos que

$$
    h4 = 2778.7275470340987 kJ/kg
$$

Hay que tomar en cuenta que las unidades que resultan de la potencia ($V*I$) son $[J/s]$
y por lo tanto tenemos que usar estas en la entalpía 4.

$$
    h_4 = 2778720\left[\frac{J}{kg}\right]
$$

Sustituyendo nuestros valores en la ecuación \ref{eq:calorimetro_electrico}

$$
    h_2 = 2778720\left[\frac{kJ}{kg}\right] - \frac{106[V]\cdot 15.5[A]}{0.01833[kg/s]} = 2689085.52\left[J/kg\right] = 2689.085\left[kJ/kg\right]
$$

Finalmente, para encontrar la calidad se reutiliza la ecuación \ref{eq:calidad}
$$
    x = \frac{2689.085\left[kJ/kg\right] - h_{f_2}}{h_{fg_2}}
$$

Las variables restantes ya se calcularon en la sección \ref{stateII}

$$
    x = \frac{2689.085\left[\frac{kJ}{kg}\right] - 698.53\left[\frac{kJ}{kg}\right]}
                {2763.12\left[\frac{kJ}{kg}\right] - 698.53\left[\frac{kJ}{kg}\right]}
    = 0.9641
$$


% ============================
% CALCULO DE ETA
% ============================
\subsection{Obtención de la eficiencia}


% ============================
% CALCULO DEL GV
% ============================
\subsubsection{Cálculo del gasto de vapor} \label{gasto_vapor}

Para obtener el gasto de vapor se utiliza la ecuación del equivalente de vapor, 
de la cual despejamos el gasto de vapor
\begin{equation} \label{eq:gasto_vapor}
    G_v = \frac{E_v\cdot h_{fg}}{h_2 - h_1}
\end{equation}

Como dato se tiene
$$E_v = 1050\left[\frac{kg}{h}\right]$$
$$h_{fg} = 2274.67\left[\frac{kJ}{kg}\right]$$

Sustituyendo en la ecuación \ref{eq:gasto_vapor}

$$
    G_v = \frac{1050\left[\frac{kg}{h}\right] \cdot 2274.67\left[\frac{kJ}{kg}\right] }
    {(h_2 - h_1)\left[\frac{kJ}{kg}\right]}
$$

Para la entalpía se utilizará la calculada a partir de las mediciones tomadas
del calorímetro eléctrico, obtenida en la sección \ref{x_calorimetro_electrico}. 
La entalpía 1 fue obtenida en la sección \ref{entalpia_I}

$$
    G_v = \frac{1050\left[\frac{kg}{h}\right] \cdot 2274.67\left[\frac{kJ}{kg}\right] }
    {(2689,085 - 133,952)\left[\frac{kJ}{kg}\right]} = 934.74\left[\frac{kg}{h}\right]
$$

Para dejarlo en término de segundos
$$
    G_v = 
    934.74\left[\frac{kg}{h}\right] \cdot \frac{1[h]}{60[min]}\cdot \frac{1[min]}{60[s]} 
    = 0.25965\left[\frac{kg}{s}\right]
$$

% ============================
% CALCULO DEL GC
% ============================
\subsubsection{Cálculo del gasto de combustible}

Para el gasto de combustible se utiliza la ecuación

\begin{equation}
    G_c = \frac{a\cdot a\cdot A}{t}\times\rho_{comb}
\end{equation}

Conociendo que $\rho_{comb} = 800\left[\frac{kg}{m^3}\right]$, y de los datos de tablas

\begin{tabularx}{0.8\textwidth} { 
  | >{\raggedright\arraybackslash}X 
  | >{\centering\arraybackslash}X 
  | >{\raggedleft\arraybackslash}X | }
 \hline
 $a$ & $30[cm]$  \\
\hline
 $A$ & $5[cm]$  \\
\hline
 $t$ & $2[min], 57[s]$  \\
\hline
\end{tabularx}

Convirtiendo las unidades al SI

\begin{tabularx}{0.8\textwidth} { 
  | >{\raggedright\arraybackslash}X 
  | >{\centering\arraybackslash}X 
  | >{\raggedleft\arraybackslash}X | }
 \hline
 $a$ & $0.30[m]$  \\
\hline
 $A$ & $0.05[m]$  \\
\hline
 $t$ & $177 [s]$  \\
\hline
\end{tabularx}

$$
    G_c = \frac{0.30[m]\cdot 0.30[m] \cdot 0.05[m]}{177[s]}\times 800\left[\frac{kg}{m^3}\right]
$$

$$
    = 0.020338\left[\frac{kg}{s}\right]
$$

\subsubsection{Cálculo de la eficiencia}

Finalmente sustituimos los valores anteriores en la ecuación
\begin{equation}
    \eta = \frac{G_v(h_2 - h_1)}{G_c(PCA)}
\end{equation}

$$
    \eta = \frac{
        0.25965\left[\frac{kg}{s}\right](h_2 - h_1)
    }{
        0.020338\left[\frac{kg}{s}\right](PCA)
    }
$$

Conociendo $PCA = 42094\left[\frac{kJ}{kg}\right]$ y utilizando los valores de entalpía 
calculados en las secciones \ref{entalpia_I} y \ref{x_calorimetro_electrico}

$$
    \eta = \frac{
        0.25965\left[\frac{kg}{s}\right](2689,085 - 133,952)\left[\frac{kJ}{kg}\right]
    }{
        0.020338\left[\frac{kg}{s}\right](42094)\left[\frac{kJ}{kg}\right]
    } = 0.774949
$$

\subsection{Obtención de la capacidad nominal y la capacidad real}

Para obtener estas capacidades se utilizan los datos

\begin{tabularx}{0.8\textwidth} { 
  | >{\raggedright\arraybackslash}X 
  | >{\centering\arraybackslash}X 
  | >{\raggedleft\arraybackslash}X | }
 \hline
 Sup.Calef & $46.61[m^2]$  \\
\hline
\end{tabularx}

Sustituyendo en la ecuación

\begin{equation}
    C_N = \frac{{Sup. Calef}}{0.93\left[\frac{m^2}{CC}\right]}\times 8510\left[\frac{kcal}{h}\right]
\end{equation}

obtenemos

$$C_N = \frac{
            46.61[m^2]
        }{
            0.93\left[\frac{m^2}{CC}\right]
        }
        \times 8510\left[\frac{kcal}{h}\right] = 426506.55\left[\frac{kcal}{h}\right]
$$

Para convertirlo a $[kJ/s]$ 
$$
    C_N = 426506.55\left[\frac{kcal}{h}\right]
    \times \frac{4.186[kJ]}{1[kcal]} \times \frac{1[h]}{3600[s]} = 495.93\left[\frac{kJ}{s}\right]
$$


Finalmente, la capacidad real se obtiene por medio de
\begin{equation}
    C_R = G_v(h_2 - h_1)
\end{equation}

Y utilizando los datos de las secciones \ref{entalpia_I}, \ref{x_calorimetro_electrico} y el
gasto de vapor calculado en la sección anterior

$$
    C_R = 0.25965\left[\frac{kg}{s}\right]
    (2689.085 - 133.952)\left[\frac{kJ}{kg}\right] = 638.78325 \left[\frac{kJ}{s}\right]
$$

\subsection{Factor de sobrecarga y factor de vaporización}

El factor de sobrecarga se obtiene por medio de la siguiente ecuación

\begin{equation}
    F_{sc} = \frac{C_R}{C_N} 
\end{equation}

y sustituyendo

$$
    F_{sc} = \frac{638.78325 \left[\frac{kJ}{s}\right]}
    {495.93\left[\frac{kJ}{s}\right]} = 1.2880
$$

Por último, el factor de vaporización utiliza la ecuación

\begin{equation}
    F_v = \frac{E_v}{G_v}
\end{equation}

Sustituimos el gasto de vapor y el equivalente de vaporización revisados en la sección \ref{gasto_vapor}

$$
    F_v = \frac{1050\left[\frac{kg}{h}\right]}{934.74\left[\frac{kg}{h}\right]} 
    = 1.123307
$$


\pagebreak

% ============================
% TABLA DE RESULTADOS
% ============================
\section{Tabla de resultados}

\begin{table}[H] \label{tab:results}
\centering
\renewcommand{\arraystretch}{1.4}
    \begin{tabular}{|l|l|}
    \hline
    \textbf{Variable} & \textbf{Resultado} \\
    $x_1$ & $0.97$ \\
    \hline
    $x_2$ & $0.96$ \\
    \hline
    $\eta$ & $77.49\%$ \\
    \hline
    $C_N$ & $495.93\left[\frac{kJ}{s}\right]$ \\
    \hline
    $C_R$ & $638.78\left[\frac{kJ}{s}\right]$ \\
    \hline
    $F_{sc}$ & $1.28$ \\
    \hline
    $F_{v}$ & $1.12$ \\
    \hline
    \end{tabular}
\caption{Resultados obtenidos.}
\end{table}

% ============================
% ANÁLISIS, CONCLUSIONES Y BIBLIOGRAFÍA
% ============================
\section{Análisis de resultados}

A partir de los primeros dos cálculos de la calidad del vapor de la caldera se obtiene
una diferencia aproximada de $0.014411$. Esta es una diferencia despreciable y podemos observar que ambos
calorímetros convergen a un valor aun al ser diferentes métodos de medición.\linebreak

La eficiencia, de acuerdo a 
\href{https://coalbiomassboiler.com/industrial-steam-boiler-efficiency/}{la eficiencia de una caldera industrial},
se acerca bastante al estándar de $80\%$. 
Los errores de medición son aceptables pero pueden eliminarse, como la lectura de las temperaturas
en los termómetros, que se utilizaron para el cálculo de la entalpía 1 y para entrar a
tablas, así como el tiempo aproximado tomado con el cronómetro, contribuyen a este error. 
Además, no tenemos información sobre el mantenimiento que recibe el equipo.\linebreak

Por otro lado, la capacidad real supera por mucho a la capacidad nominal, 
lo que indica otro error en las mediciones. El factor de sobrecarga muestra un incremento del
$28\%$ con respecto a la capacidad nominal. Esto significa que la 
caldera está operando por encima de su punto de diseño. Con respecto al factor de vaporización, 
se muestra una proporción
a favor del equivalente de vaporización del $12\%$.

\section{Conclusiones}

Se determinó la calidad del vapor de la caldera por dos métodos independientes:
el calorímetro de estrangulamiento ($x_1 = 0.9786$) y el calorímetro eléctrico
($x_2 = 0.9641$). La diferencia entre ambos (1.5\%) es atribuible a las pérdidas
de calor en la resistencia eléctrica y a la suposición de proceso isentálpico
en el estrangulamiento. La calidad promedio ($x = 0.971$) indica un vapor con
aproximadamente 3\% de humedad, aceptable para aplicaciones de calefacción pero
insuficiente para turbinas de vapor.

La eficiencia calculada ($\eta = 77.49\%$) es consistente con el rango típico
de calderas de laboratorio (70--80\%), aunque ligeramente por debajo del estándar
industrial de 80\%. Las principales fuentes de error son la lectura de temperaturas
en los termómetros ($\pm 1^\circ$C), el tiempo medido con cronómetro ($\pm 0.5$ s)
y la posible pérdida de calor en el calorímetro eléctrico.

El factor de sobrecarga ($F_{sc} = 1.28$) indica que la caldera opera 28\% por
encima de su capacidad nominal, lo que podría deberse a una sobreestimación del
gasto de vapor o a una subestimación de la capacidad nominal. El factor de
vaporización ($F_v = 1.12$) es cercano a 1, lo que indica que el
equivalente de vaporización real es 12\% mayor que el nominal.

\section{Bibliografía}

\begin{itemize}
    \item Bell, I. H., Wronski, J., Quoilin, S., \& Lemort, V. (2014). 
    Pure and pseudo-pure fluid thermophysical property evaluation and 
    the open-source thermophysical property library CoolProp. 
    \textit{Industrial \& Engineering Chemistry Research}, 53(6), 
    2498--2508. \url{https://doi.org/10.1021/ie4033999}
    
    \item Cengel, Y. A., \& Boles, M. A. (2015). 
    \textit{Termodinámica} (8a ed.). McGraw-Hill Education.
    
    \item Coalmass Boiler. (2024). \textit{Industrial Steam Boiler Efficiency}. 
    Recuperado de 
    \url{https://coalbiomassboiler.com/industrial-steam-boiler-efficiency/}
    
    \item López Maldonado, S. (2026). \textit{Hoja de calculos calderas.ppt} [Material de curso]. 
    Facultad de Ingeniería, Universidad Nacional Autónoma de México.
    
    \item Python Software Foundation. (2024). \textit{Python Language 
    Reference, version 3.12}. \url{https://docs.python.org/3/}
\end{itemize}

\end{document}